\chapter{Experimentos e inferencia estadística}\label{ch:experiments-and-inference}

% Alcance: contraste de hipótesis nula, potencia estadística y tamaño de muestra, regresión OLS,
%   diseño de pruebas A/B para decisiones de negocio. Capítulo nuevo (cierra la brecha de diseño experimental).
% El libro emplea pruebas A/B y RCT de incrementalidad (\cref{ch:attribution-and-measurement})
%   sin enseñar la maquinaria subyacente; este capítulo la provee, ampliando \cref{ch:decisions-under-uncertainty}.
% Redacción según el estilo de la obra: prosa continua + recuadros example/admonition; ejemplos con la SaaS de referencia.

\section{El marco de la hipótesis nula}\label{sec:null-hypothesis}
% complexity: 2

Toda decisión de producto es una apuesta. El marco de la hipótesis nula obliga al tomador de
decisiones a cuantificar cuán improbables serían los datos observados si la intervención no
tuviera \emph{ningún} efecto, y luego decidir si esa improbabilidad es suficientemente baja
para actuar. La intuición gerencial es simple: recolectar evidencia, preguntar con qué
frecuencia datos tan extremos aparecerían por azar bajo el statu quo, y rechazar el statu
quo solo cuando esa probabilidad cae por debajo de un umbral acordado con antelación. La
disciplina radica en el compromiso previo: elegir el umbral antes de ver los resultados, no
después.

Bajo la hipótesis nula \(H_0: \mu_1 = \mu_2\) de que dos grupos tienen la misma media, el
estadístico de prueba para muestras grandes es la distancia estandarizada entre las medias
grupales observadas, escalada por el error estándar de su diferencia:
\begin{equation}\label{eq:z-test}
  z \;=\; \frac{\bar{x}_1 - \bar{x}_2}%
               {\sqrt{\dfrac{\sigma_1^2}{n_1}+\dfrac{\sigma_2^2}{n_2}}}.
\end{equation}
Se rechaza \(H_0\) cuando \(|z| > z_{\alpha/2}\), el valor crítico para una prueba bilateral
al nivel de significancia \(\alpha\). Para el convencional \(\alpha=0.05\), \(z_{\alpha/2}=1.96\).
Para resultados de proporción---tasas de conversión, tasas de conversión de prueba a pago---\(\sigma_i^2 = p_i(1-p_i)\).
La fórmula se apoya en el Teorema del Límite Central, que requiere muestras suficientemente
grandes para que la distribución muestral de la media sea aproximadamente normal; en la
práctica, \(n \geq 30\) por grupo es un piso habitual, aunque para proporciones pequeñas se
necesitan muestras mucho mayores antes de que la aproximación sea confiable.

\begin{example}[Prueba A/B del \emph{onboarding} SaaS]
El equipo de producto prueba un flujo de \emph{onboarding} rediseñado. Control: \num{200}
pruebas, \num{50} convierten a pago (\(\hat{p}_c = \qty{25}{\percent}\)). Tratamiento:
\num{200} pruebas, \num{62} convierten (\(\hat{p}_t = \qty{31}{\percent}\)). La decisión:
¿el aumento de \qty{6}{\percent} es real o es ruido?

Se fija \(\alpha=0.05\), por lo que el valor crítico es \(z_{0.025}=1.96\). Se estima la
proporción agrupada bajo \(H_0\) como \(\hat{p} = (50+62)/(200+200) = 0.28\), lo que da
\(\hat{\sigma}^2 = 0.28 \times 0.72 = 0.2016\). El error estándar es
\[
  \mathrm{SE} \;=\; \sqrt{\frac{0.2016}{200}+\frac{0.2016}{200}}
               \;=\; \sqrt{0.002016}
               \;\approx\; 0.0449.
\]
Aplicando \cref{eq:z-test},
\[
  z \;=\; \frac{0.31 - 0.25}{0.0449} \;\approx\; 1.34.
\]
Como \(|z|=1.34 < 1.96\), el resultado \emph{no} es estadísticamente significativo con
\(n=200\) por brazo. El aumento de \qty{6}{\percent} bien podría deberse al azar. La
implicación no es que el tratamiento falle; es que el experimento es demasiado pequeño para
distinguir un efecto real de esta magnitud del ruido muestral. Recolectar más datos es el
único remedio---cuántos más es el tema de \cref{sec:power-and-sample-size}.
\end{example}

\begin{admonition}{Advertencia}
El valor p es la probabilidad de observar un estadístico de prueba al menos tan extremo
\emph{si la hipótesis nula fuera verdadera}. No es, en absoluto, la probabilidad de que la
hipótesis nula sea verdadera, ni de que la alternativa sea falsa. Un valor p de \num{0.03}
no significa que haya un \qty{97}{\percent} de probabilidad de que el tratamiento funcione;
significa que los datos son moderadamente incompatibles con «ningún efecto». Igualmente, la
significancia estadística no dice nada sobre el tamaño del efecto: un aumento trivial de
\qty{0.1}{\percent} alcanzará \(p<0.05\) con una muestra suficientemente grande. Las
decisiones de negocio requieren tanto un resultado significativo \emph{como} un efecto lo
bastante grande para importar económicamente.
\end{admonition}

El estadístico de prueba en \cref{eq:z-test} proporciona la maquinaria; \cref{sec:power-and-sample-size}
establece cuántas observaciones se necesitan antes de que esa maquinaria sea lo bastante
potente para merecer ejecutarse. \textcite{anderson2019statistics} desarrollan la teoría
completa de la distribución muestral; \textcite{wooldridge2020introductory} la extiende a
los escenarios de regresión donde múltiples covariables entran simultáneamente.


\section{Potencia estadística y tamaño de muestra}\label{sec:power-and-sample-size}
% complexity: 4

Un experimento incapaz de detectar el efecto para el que fue diseñado no es un resultado
neutro; es un desperdicio. La potencia estadística---la probabilidad de rechazar correctamente
una hipótesis nula falsa---cuantifica el riesgo de ese fracaso. Una potencia baja significa
que el experimento devolverá con más frecuencia un resultado no concluyente incluso cuando
el tratamiento funciona genuinamente, generando una cadena de hallazgos de «diferencia no
significativa» que se interpretan erróneamente como evidencia de que el producto está bien.
El costo de negocio es simétrico: los experimentos con baja potencia desperdician tiempo en
mejoras reales y dan falsa tranquilidad cuando un cambio está causando daño.

La potencia depende de tres insumos que el experimentador controla: el tamaño de muestra
\(n\), el nivel de significancia \(\alpha\) y el efecto mínimo detectable \(\delta\). Para
una prueba de dos brazos con grupos de igual tamaño y varianza común \(\sigma^2\), el
tamaño de muestra requerido por brazo es:
\begin{equation}\label{eq:sample-size}
  n \;=\; \frac{2\,(z_{\alpha/2}+z_{\beta})^{2}\,\sigma^{2}}{\delta^{2}}.
\end{equation}
Aquí \(z_{\beta} = \Phi^{-1}(1-\beta)\) es el valor crítico para la potencia deseada
\(1-\beta\); con \qty{80}{\percent} de potencia, \(z_{\beta}=0.84\). La relación de cuadrado
inverso entre \(\delta\) y \(n\) es el hecho central del diseño experimental: reducir a la
mitad el tamaño del efecto de interés \emph{cuadruplica} el tamaño de muestra requerido. La
intuición es que las señales más pequeñas son más difíciles de distinguir del ruido, por lo
que se necesitan más observaciones para ahogar ese ruido. \Cref{fig:power-curve} grafica la
potencia en función del tamaño del efecto para tres tamaños de muestra, haciendo concreta
la disyuntiva.

\begin{figure}[htbp]
  \centering
  \includegraphics[width=0.86\linewidth]{power-curve}
  \caption{Potencia estadística frente al tamaño del efecto absoluto (incremento en la tasa
    de conversión sobre una base de \qty{25}{\percent}) con \(\alpha=0.05\), prueba bilateral,
    para \num{100}, \num{200} y \num{550} observaciones por brazo. La línea punteada marca
    el objetivo de \qty{80}{\percent} de potencia; el diamante marca el punto de diseño
    requerido. Detectar el incremento de \qty{6}{pp} de \cref{sec:null-hypothesis} con
    \qty{80}{\percent} de potencia requiere aproximadamente \num{879} observaciones por brazo.}
  \label{fig:power-curve}
\end{figure}

\begin{example}[Cálculo del tamaño de muestra para la prueba A/B de \emph{onboarding}]
¿Cuántas pruebas se necesitan para detectar el incremento de \qty{25}{\percent} \(\to\)
\qty{31}{\percent} en \emph{onboarding} con \qty{80}{\percent} de potencia y \(\alpha=0.05\)?
El efecto mínimo detectable es \(\delta=0.06\). Se usa la proporción agrupada
\(\hat{p}=0.28\) para estimar \(\sigma^2 = 0.28\times 0.72 = 0.2016\) y se sustituye en
\cref{eq:sample-size}:
\[
  n \;=\; \frac{2\,(1.96+0.84)^{2}\times 0.2016}{0.06^{2}}
      \;=\; \frac{2\times 7.84 \times 0.2016}{0.0036}
      \;\approx\; \num{879} \text{ por brazo}.
\]
El \emph{funnel} entrega \num{800} pruebas por mes, de modo que con \num{400} por brazo el
experimento alcanza \num{879} en aproximadamente \num{2.2} meses---un horizonte manejable.
¿Vale la pena la espera? El valor de confirmar el incremento: \qty{6}{pp} más en conversiones
de pago \(\times\) \num{200} nuevos clientes/mes \(\times \money{3150}\) CLV por cliente
\(= \money{37800}\) de contribución mensual incremental---un flujo recurrente que vale
\(\money{37800}/0.08 \approx \money{472500}\) a perpetuidad al WACC del \qty{11}{\percent}
de la empresa menos \qty{3}{\percent} de crecimiento (\cref{eq:gordon}). Frente a un retraso
de dos meses, esa razón favorece abrumadoramente realizar una prueba disciplinada antes que
adivinar y potencialmente lanzar un cambio que no hace nada o causa daño.
\end{example}

\begin{admonition}{Advertencia}
La detención opcional---revisar los resultados a diario y detener el experimento en el momento
en que \(p<0.05\)---infla el error de Tipo I real muy por encima del \qty{5}{\percent} nominal.
Si se revisan los resultados diez veces durante un experimento, la probabilidad de un resultado
significativo espurio puede superar el \qty{30}{\percent}. La solución es directa: pre-registrar
la regla de detención y el \(n\) requerido antes del lanzamiento, y leer los resultados exactamente
una vez. Las plataformas que ofrecen pruebas secuenciales «siempre válidas» (\textcite{johari2022always})
permiten monitoreo continuo sin inflar el error, pero requieren un estadístico de prueba
subyacente diferente (una prueba de razón de probabilidad secuencial mixta, mSPRT) y típicamente
necesitan muestras más grandes que la fórmula de horizonte fijo anterior para alcanzar la misma
potencia.
\end{admonition}

El cálculo del tamaño de muestra en \cref{eq:sample-size} es el compromiso previo que hace
honesto a \cref{eq:z-test}. Sin él, un equipo siempre puede encontrar significancia ejecutando
el experimento por más tiempo---una práctica que destruye la validez estadística de todo lo
que toca.


\section{Regresión: control de covariables}\label{sec:regression}
% complexity: 3

Los experimentos aleatorizados establecen causalidad por construcción: la asignación aleatoria
equilibra todos los confusores observados y no observados entre tratamiento y control, por lo
que cualquier diferencia en los resultados debe reflejar el tratamiento. Pero muchas preguntas
de negocio no pueden aleatorizarse: el \emph{churn} depende de la antigüedad, la profundidad
del producto, el \emph{tier} del plan y el canal de ventas, y los clientes no fueron asignados
aleatoriamente a esos estados. La regresión de mínimos cuadrados ordinarios (OLS) responde a
la pregunta «¿cuánto importa cada factor, manteniendo fijos todos los demás?» Es la herramienta
central del análisis de negocio observacional precisamente porque es el instrumento más sencillo
para mantener algunas cosas constantes mientras se examinan otras.

El modelo de regresión lineal múltiple postula que el resultado \(y_i\) es una función lineal
de \(k\) variables explicativas más un término de error que captura todo lo que el modelo
omite:
\begin{equation}\label{eq:ols}
  y_i \;=\; \beta_0 + \beta_1 x_{1i} + \cdots + \beta_k x_{ki} + \varepsilon_i,
  \qquad
  \hat{\boldsymbol{\beta}} \;=\; (\mathbf{X}^{\!\top}\mathbf{X})^{-1}\mathbf{X}^{\!\top}\mathbf{y}.
\end{equation}
La solución matricial \(\hat{\boldsymbol{\beta}}\) minimiza la suma de cuadrados de los
residuos. Por el teorema de Gauss--Markov, OLS es el Mejor Estimador Lineal Insesgado (BLUE,
por sus siglas en inglés) cuando los errores tienen media cero, varianza constante
(homocedasticidad) y no están correlacionados con los regresores. Cada coeficiente
\(\hat{\beta}_j\) mide el cambio marginal en \(y\) ante un cambio unitario en \(x_j\),
\emph{ceteris paribus}. Esa condición de «todo lo demás igual» es la fuente tanto del poder
del método como de su peligro: mantiene constante solo lo que está en el modelo, y los
confusores no observados absorbidos en \(\varepsilon\) sesgará las estimaciones si se
correlacionan con los regresores.

\begin{example}[Cuantificación del \emph{moat} de costo de cambio por integraciones]
El equipo de éxito del cliente quiere saber cuánto protege contra el \emph{churn} la
profundidad del producto. Regresa la tasa mensual de \emph{churn} sobre tres variables a
nivel de cliente: antigüedad en meses, número de integraciones activas y \emph{tier} del
plan (1~=~inicial, 2~=~profesional, 3~=~empresarial). El modelo estimado sobre la cohorte
de \num{1200} clientes de la empresa devuelve:
\begin{align*}
  \widehat{\text{churn}}_i \;=\;{}
    & 0.032
      - 0.00012\,\text{tenure}_i
      - 0.0080\,\text{integrations}_i
      - 0.0031\,\text{tier}_i.
\end{align*}
El coeficiente sobre integraciones es \(\hat{\beta}_{\text{int}} = -0.008\): cada integración
activa adicional reduce el \emph{churn} mensual en \qty{0.8}{pp}, manteniendo antigüedad y
\emph{tier} fijos. Con el \emph{churn} mensual base de la empresa en \qty{0.65}{\percent},
un cliente con cinco integraciones en lugar de una tiene un \emph{churn} previsto de
\(0.65\% - 4\times0.8\% = 0.65\% - 3.2\% \approx \) prácticamente cero---un perfil de
retención casi perfecto. Este coeficiente cuantifica el \emph{moat} de costo de cambio
descrito en \cref{sec:moat-taxonomy}: cada integración que el cliente construye es un hilo
que lo vincula a la plataforma. La implicación para la inversión en producto: una
funcionalidad que lleve a los clientes a añadir una integración más vale, en CLV anual
esperado, la reducción en \emph{churn} multiplicada por el CLV por cliente de \money{3150}.
\end{example}

\begin{admonition}{Advertencia}
Los coeficientes de regresión son correlacionales, no causales, a menos que el mecanismo de
asignación esté aleatorizado o se use un instrumento válido. La correlación negativa entre
integraciones y \emph{churn} podría reflejar en parte \emph{selección}: los clientes que
construyen integraciones profundas pueden haber elegido el producto porque planeaban quedarse,
no solo porque las integraciones elevan el costo de cambio. La única manera de aislar la
contribución causal de las integraciones es un experimento aleatorizado---incentivar a un
subconjunto aleatorio de usuarios a añadir integraciones adicionales y comparar el
\emph{churn} posterior (\cref{sec:ab-test-design}). La regresión observacional controla los
confusores medidos; la aleatorización controla todo.
\end{admonition}

\textcite{wooldridge2020introductory} demuestra el resultado de Gauss--Markov y detalla los
diagnósticos necesarios para evaluar si sus supuestos se cumplen; \textcite{angrist2009mostly}
conecta OLS con la inferencia causal y explica cuándo se necesitan variables instrumentales
o diseños de diferencias en diferencias. El coeficiente de costo de cambio anterior es un
ejemplo de una clase más amplia de mecanismos de creación de valor catalogados en
\cref{sec:moat-taxonomy}; el mismo \emph{framework} de regresión aplicado a datos a nivel
geográfico subyace a la medición de incrementalidad en \cref{sec:incrementality}.


\section{Diseño de pruebas A/B para decisiones de negocio}\label{sec:ab-test-design}
% complexity: 3

Conocer la mecánica de \cref{eq:z-test} y \cref{eq:sample-size} es necesario pero no
suficiente. Una prueba bien diseñada requiere cuatro decisiones previas que nada tienen que
ver con la fórmula: qué métrica primaria optimizar, qué métricas de protección no pueden
degradarse, a qué nivel aleatorizar los usuarios y por cuánto tiempo ejecutar el experimento.
Omitir cualquiera de ellas convierte un cálculo estadísticamente válido en un experimento
comercialmente poco confiable. La disciplina del practicante es anotar las cuatro antes de
que se recolecte cualquier dato.

La \textbf{métrica primaria} (criterio de evaluación global, OEC) debe ser la métrica más
directamente vinculada al valor a largo plazo. Para una prueba de \emph{onboarding} SaaS,
la conversión de prueba a pago es la métrica primaria correcta: determina directamente los
\num{200} nuevos clientes por mes de la hoja de datos y, a través del CLV, el incremento de
\money{37800}/mes calculado en \cref{sec:power-and-sample-size}. El ingreso por prueba es
tentador pero demasiado ruidoso; la sesión a prueba es anterior y puede moverse sin beneficio
corriente abajo. Las \textbf{métricas de protección} salvaguardan la salud básica del negocio:
si el tratamiento mejora la conversión de prueba a pago pero ralentiza la carga de página en
dos segundos, la ganancia a corto plazo quedará borrada por el menor compromiso a largo plazo.
Preespecificar las métricas de protección evita que el equipo lance un tratamiento que gana
en la métrica primaria mientras daña silenciosamente el \emph{funnel}.

La \textbf{unidad de aleatorización} determina qué violaciones de las métricas de protección
son detectables y rige si se cumple el Supuesto del Valor de Tratamiento de Unidad Estable
(SUTVA). Aleatorizar a nivel de sesión es sencillo, pero puede contaminar los resultados si
el mismo usuario cae en ambos brazos a lo largo de sesiones. Para una prueba de
\emph{onboarding}, la asignación a nivel de usuario o dispositivo es la elección correcta:
el tratamiento es la experiencia completa de \emph{onboarding}, no una sola visita de página,
y el evento de conversión ocurre a lo largo de múltiples sesiones. La aleatorización a nivel
de sesión mezclaría experiencias de control y tratamiento para el mismo usuario en prueba,
sesgando la tasa de conversión observada en ambos brazos. Cuando los efectos de red
significan que el resultado de un usuario depende del estado de tratamiento de otro---habitual
en productos sociales y mercados de dos lados---la aleatorización a nivel geográfico (como en
el \emph{framework} de incrementalidad de \cref{sec:incrementality} y el estimador iROAS en
\cref{eq:iroas}) es el único diseño que evita las violaciones del SUTVA.

\begin{example}[Diseño de la prueba de la página de precios]
El equipo de \emph{growth} quiere probar un rediseño de la página de precios que hace más
prominente el plan profesional. El \emph{framework} de decisión es:

\begin{description}
  \item[Métrica primaria] Tasa de conversión de prueba a pago. Base actual \qty{25}{\percent};
    efecto mínimo detectable \qty{4}{pp} (un incremento de \qty{25}{\percent} a
    \qty{29}{\percent}), anclado al margen esperado de \num{8} conversiones de pago
    incrementales por mes (\(0.04\times 200\)) multiplicadas por \money{3150} CLV
    \(= \money{25200}\)/mes.
  \item[Métricas de protección] (a) Tasa de sesión a prueba: no debe caer más de \qty{2}{pp}
    (una caída relativa del \qty{10}{\percent} eliminaría la ganancia en el \emph{funnel});
    (b) Volumen de tickets de soporte: no más del \qty{15}{\percent} de aumento (una página
    de precios muy interrogada infla el costo de atención).
  \item[Unidad de aleatorización] Usuario (\emph{cookie} + inicio de sesión), asignado en
    la primera sesión. Un usuario en el brazo de tratamiento ve la página de precios
    rediseñada en cada visita; un usuario de control siempre ve la original.
  \item[Tamaño de muestra y duración] Usando \cref{eq:sample-size} con \(\delta=0.04\),
    \(\sigma^2=0.25\times 0.75=0.1875\), \(\alpha=0.05\), \(1-\beta=0.80\):
    \[
      n \;=\; \frac{2\,(1.96+0.84)^{2}\times 0.1875}{0.04^{2}} \approx \num{1838}
        \text{ por brazo.}
    \]
    Con \num{800} pruebas por mes divididas equitativamente entre dos brazos (del \emph{funnel}
    en la hoja de datos), la prueba requiere aproximadamente \num{4.6} meses de datos---una
    espera significativa. El equipo debe o bien aceptar un efecto mínimo detectable mayor
    (digamos, \qty{6}{pp} \(\to\) \num{879} por brazo \(\approx\) 2.2 meses) o invertir en
    reducir la varianza de conversión mediante ajuste de covariables previo al experimento al
    estilo CUPED. Pre-registrar la regla de detención a tres semanas con \num{879}-por-brazo
    y leer los resultados exactamente una vez elimina el sesgo de detención opcional.
\end{description}

La implicación: la restricción vinculante en la velocidad de experimentación no es la teoría
estadística sino el tamaño del \emph{funnel}. Una empresa que adquiere solo \num{200} clientes
de pago por mes necesita meses para ejecutar pruebas con potencia razonable sobre conversiones
corriente abajo. Invertir en la parte superior del \emph{funnel} de adquisición
(\cref{ch:decisions-under-uncertainty}) amplía la superficie experimental tanto como amplía
los ingresos.
\end{example}

\begin{admonition}{Nota}
Pre-registrar el diseño---métrica primaria, métricas de protección, unidad de aleatorización,
tamaño de muestra y regla de detención---antes del lanzamiento no es burocracia; es el
mecanismo que mantiene honesto el valor p del análisis. Sin pre-registro, un equipo siempre
puede encontrar una métrica que se movió, una ventana temporal que muestra significancia o
un subgrupo que confirma su hipótesis previa. \textcite{kohavi2020trustworthy} documentan
los modos de fallo más comunes de los programas de prueba A/B a gran escala en Microsoft y
otras empresas; \textcite{gordon2019comparison} muestran que los métodos observacionales
pueden perder efectos causales reales por márgenes grandes e inconsistentes, lo que es
precisamente la razón por la que una prueba correctamente diseñada---incluso una lenta---vale
la pena ejecutar cuando los ingresos en juego justifican el cálculo del iROAS en \cref{eq:iroas}.
\end{admonition}

El \emph{framework} de diseño de pruebas aquí presentado es el complemento operativo de la
caja de herramientas de incrementalidad en \cref{sec:incrementality}: ambos responden la
misma pregunta causal---¿produce esta intervención más ingresos que el contrafactual?---pero
a distintas granularidades. Los experimentos de producto aleatorizan usuarios dentro de una
superficie de producto y miden la conversión; los experimentos de \emph{holdout} geográfico
aleatorizan mercados geográficos y miden el iROAS de pleno \emph{funnel} del gasto en
medios. \textcite{johnson2017ghost} demuestran la magnitud de la sobreestimación del anunciante
cuando la aleatorización a nivel geográfico se omite en favor de la atribución reportada por
la propia plataforma. El hilo conductor entre los tres capítulos---este capítulo,
\cref{ch:attribution-and-measurement} y \cref{ch:decisions-under-uncertainty}---es el marco
de resultados potenciales: un efecto causal solo existe como comparación frente a un
contrafactual que nunca se observó, y la única manera confiable de construir ese contrafactual
es la aleatorización.
